% LaTeX draft: Unit of Caring
\documentclass[11pt]{article}
\usepackage[utf8]{inputenc}
\usepackage[T1]{fontenc}
\usepackage{lmodern}
\usepackage{geometry}
\geometry{margin=1in}
\usepackage{amsmath,amssymb,amsthm,mathtools}
\usepackage{microtype}
\usepackage[numbers]{natbib}
\usepackage{hyperref}
\hypersetup{hidelinks}

\title{A Unit of Caring\\Integrity Pressure, Suffering, and Cross-Scale Aggregation}
\author{Gunnar Zarncke, aintelope}
\date{May 20, 2026}

\begin{document}
\maketitle

\begin{abstract}
We propose a formal scaffold for talking about ``how much is at stake'' in situations that involve harm, pain, or suffering. The central quantity is a scalar ``unit of caring'' that is intended to be neutral with respect to particular ethical theories but precise enough to connect to control-theoretic and information-theoretic descriptions of agents.\footnote{This is a conceptual and mathematical framing, not a completed moral theory. Many normative choices are intentionally left underspecified.}

We proceed in three steps. First, rather than grounding harm in agents' explicit preferences, we introduce an integrity-based quantity $H_i(t)$, the \emph{integrity pressure} on an agent $i$ at time $t$. Intuitively, $H_i(t)$ measures how hard the world is currently pushing against the agent's continued existence, organisation, and characteristic functioning. We treat $H_i$ as a black box that a more detailed physical or variational theory could make precise.

Second, we distinguish $H_i$ from the \emph{perceived} integrity pressure $\tilde H_i$ that the agent's controllers actually act on, and from the phenomenological quantity $S_i(t)$ that we call suffering. 
We treat suffering as a gain-modulated function of pain signals: $S_i(t) \approx P_i(t)\,G_i(\dots)$, where $P_i$ aggregates low-level drive signals, and $G_i$ is a gain determined by architectural parameters.

Third, we define a global caring functional
\[
  \mathcal L = \sum_{i} W_i D_i, \qquad D_i = \int H_i(t)\,dt,
\]
with scale- and architecture-dependent weights $W_i$. Instead of constructing $W_i$ from a specified base level of constituents, we state qualitative constraints that $W$ should satisfy (non-negativity, sublinear scaling with size, approximate invariance under coarse-graining, and bounded overcounting). Within those constraints, different choices of $W$ correspond to different---but structurally comparable---ethical views about how to trade off harms across species and scales.

We use this scaffold to define a reference ``unit of caring'' as the change in $\mathcal L$ obtained by preventing a canonical episode of integrity pressure in a reference agent, and we show how suffering, architectural parameters, and cross-scale aggregation relate to that unit.
\end{abstract}

\section{Introduction}

Many ethical and practical decisions implicitly ask the same quantitative question: how much is at stake for some agent or set of agents? Clinical pain management asks it when comparing states that differ in intensity, duration, and distress. Public health asks it when aggregating burden across populations and time. Animal welfare and AI welfare debates ask it when comparing systems that may differ not only in what happens to them, but in how much of it is globally integrated, self-indexed, or reflectively elaborated. Yet the relevant literatures are oddly fragmented. Philosophical work on value aggregation often remains abstract and weakly tied to empirical models of control and phenomenology; neuroscience of pain and metacognition typically models local mechanisms without a general cross-scale caring functional; and applied health metrics such as QALYs and DALYs provide useful aggregation schemes but collapse many architectural differences into a single burden score \cite{whitehead2010qaly,murray1996globalburden,murray2013gbd,who_daly}.

This paper proposes a formal scaffold for bridging these domains. The central object is a \emph{unit of caring}: a reference amount of ethically relevant load, defined in terms of an agent’s integrated \emph{integrity pressure}. We distinguish three levels that are often conflated. First, there is an underlying quantity $H_i(t)$, the integrity pressure on agent $i$ at time $t$, intended to capture how strongly the world is pushing that system away from continued existence, organization, and characteristic functioning. Second, there are low-level control signals, generated by local homeostatic or threat-sensitive loops, which need not map transparently onto $H_i$. Third, there is phenomenology, in particular suffering, which is not merely the presence of nociceptive or aversive signals, but their globally integrated, recursively elaborated, and self-indexed uptake.

The motivation for separating these levels is empirical as well as conceptual. Pain science has long shown that nociception, pain, distress, and disability are dissociable \cite{melzack1965gate,melzack1975mpq,melzack2001neuromatrix}. Clinical and lesion cases further show that the affective and motivational significance of pain can come apart from its sensory-discriminative aspects, and that introspective access to one’s own states is systematically lossy rather than transparent \cite{schooler2002rerepresenting,fleming2014measure,maniscalco2012metad,gerrans2024pain}. This suggests that “how bad things are for a system” cannot be read directly from self-report, from local nociceptive activation, or from a single scalar measure of intensity. At the same time, we do need some common currency if we want to compare interventions, species, architectures, or social scales.

Our proposal is deliberately intermediate in ambition. We do not claim to derive value from physics, nor to settle population ethics, nor to provide a full theory of consciousness. Instead, we introduce a family of variables that make the comparison problem more explicit. In particular, we treat suffering as a gain-modulated transformation of lower-level aversive channels, where the gain depends on architectural variables related to bandwidth, recursion depth, self-indexing, and introspective opacity. This decomposition is continuous with the metacognitive parameterization developed elsewhere in this project, where bandwidth $B$, recursion depth $d$, opacity $\tau$, and bottlenecked selfhood $\beta$ are treated as partially separable dimensions of phenomenology and self-modeling \cite{zarncke2026rainbow}.

The resulting picture is neither purely hedonistic nor purely functionalist. It is hedonically informed, in that first-person badness matters and pain-like signals are one of our clearest entry points into the problem. But it is also functionalist and control-theoretic, in that what ultimately matters is not any one report, signal, or felt intensity alone, but the relation between integrity pressure, the controllers that represent it, and the architectures that amplify or damp it. This lets us say, for example, why anesthesia, pain asymbolia \cite{gerrans2024pain}, panic, contemplative dereification, or lesion-induced changes in self-model may sharply alter suffering without straightforwardly changing underlying integrity pressure. It also lets us ask how a global caring functional should aggregate across entities of different scale without either collapsing everything into one undifferentiated total or proliferating arbitrary double counting.

The rest of the paper develops this scaffold. We begin by situating it with respect to prior work on units of pleasure and suffering, empirical pain measurement, metacognitive access, and burden aggregation. We then define integrity pressure, low-level valenced channels, and suffering as a gain-modulated global controller. Finally, we introduce a caring functional over integrated harm loads with scale-sensitive weights, while leaving the exact construction of those weights intentionally open.

\section{Related Work}

\subsection{From hedonic units to empirical burden metrics}

The aspiration to quantify morally relevant experience is old. Bentham’s hedonic calculus treated pleasure and pain as additive quantities suitable for moral arithmetic, even if the proposed dimensions---intensity, duration, certainty, fecundity, and so on---were never operationalized into a stable empirical unit \cite{bentham1789principles}. More recent informal accounts treat scarce resources, such as money spent or lives saved per dollar, as practical proxies for how much is at stake \cite{yudkowsky2009money}. Modern burden metrics such as the QALY and DALY are, in a sense, much more successful descendants of this ambition. They provide tractable aggregation schemes over time and persons, and they are now central to health economics and global burden estimation \cite{whitehead2010qaly,murray1996globalburden,murray2013gbd,who_daly}. However, these frameworks largely presuppose rather than explain the structure of the relevant burden. They tell us how to aggregate burden once a health-state valuation is available, but they do not specify how local control signals, phenomenology, self-modeling, and architecture combine to produce that burden in the first place.

Our proposal is complementary rather than competitive. It is compatible with later aggregation into QALY- or DALY-like objects, but it aims to factor the pre-aggregation problem into more mechanistically interpretable components. In particular, it distinguishes underlying integrity load from the signals by which an organism detects and responds to that load, and from the reflectively elaborated suffering that may or may not accompany those signals.

\subsection{Pain as control, representation, and multidimensional experience}

Modern pain science strongly motivates such a separation. Gate Control Theory already broke with the idea that pain is a simple readout of tissue damage, emphasizing spinal and supraspinal modulation by context, expectation, and competing inputs \cite{melzack1965gate}. The McGill Pain Questionnaire further made clear that pain is multidimensional: sensory, affective, and evaluative components can vary partly independently and require richer description than a single intensity rating \cite{melzack1975mpq}. Melzack’s later neuromatrix framework generalized this further, treating pain as an output of distributed brain processes integrating bodily, cognitive, and affective information rather than as a direct input from nociceptors \cite{melzack2001neuromatrix,wall1999pain}.

This literature supports two claims central to the present paper. First, nociception, pain, and suffering are not identical. Second, the motivational significance of a signal depends on how it is integrated into larger representational and regulatory structures. Zaki and colleagues sharpen this point by distinguishing nociceptive pain from empathic pain and identifying partly shared affective circuitry that tracks badness more generally rather than only tissue damage \cite{zaki2016anatomy}. Gerrans, in turn, offers an active allostatic inference account in which pain and suffering are mediated by self-modeling and allostatic control, making explicit the role of the self in transforming local aversive signals into experienced suffering \cite{gerrans2024pain}. This is close in spirit to our claim that low-level aversive channels are only one input to suffering, which depends additionally on global integration, recursion, self-indexing, and opacity.

\subsection{Measuring suffering and the limits of objectivity}

The empirical measurement literature is similarly instructive. Self-report remains the clinical gold standard \cite{iasp2020pain}, but it is noisy, compressed, and architecture-dependent. Attempts to create an objective unit of pain, such as early dolorimetry, were historically important but did not yield a robust universal scale \cite{hardy1948dol,wall1999pain}. Contemporary neuroimaging approaches are more sophisticated \cite{wager2013fmripain}, yet the consensus remains that brain-based measures cannot simply replace first-person report \cite{thacker2012firstperson,davis2017brainimaging}. Davis et al.\ review the medical, legal, and ethical issues around pain neuroimaging and conclude that such measures are informative but not authoritative in isolation. This fits naturally with our framework: suffering signals are evidence about underlying integrity pressure, but they are not identical to it, and their evidential value depends on calibration.

\subsection{Metacognition, opacity, and recursive self-modeling}

A second cluster of relevant work concerns metacognition and self-representation. Schooler’s distinction between consciousness and meta-consciousness is particularly useful here: one may undergo an experience without simultaneously representing oneself as undergoing it, and the transition to meta-awareness is itself intermittent and structured \cite{schooler2002rerepresenting}. Signal-detection-based measures such as meta-$d'$ and meta-$d'/d'$ provide quantitative ways of capturing how much information about first-order performance is available to higher-order monitoring \cite{maniscalco2012metad,fleming2014measure}. These measures do not by themselves solve the “hard problem,” but they do provide tractable handles on introspective opacity and higher-order access.

This matters because a plausible model of suffering must explain not only why some signals are globally available, but why they appear intrinsic, inescapable, and personally binding. The broader metacognitive program associated with the Rainbow parameterization \cite{zarncke2026rainbow} treats bandwidth, opacity, recursion depth, and self-bottleneck strength as partially separable dimensions of phenomenology and subjectivity. Complementary scalar indices of integrated brain state have also been proposed \cite{casali2013pci}. On that picture, suffering is not just “pain plus more pain,” but pain-like or aversive signals routed through a globally available, recursively elaborated, and self-indexed system under limited transparency. While the present paper does not depend on any one consciousness theory, it draws directly on this decomposition.

\subsection{Aggregation across agents and scales}

Finally, there is the aggregation problem. Health economics and burden-of-disease work provide one important template: define a per-agent or per-person burden, then sum or weight it across time and population \cite{whitehead2010qaly,murray1996globalburden,murray2013gbd}. Population well-being surveys offer an alternative template based on aggregated self-reported life satisfaction \cite{helliwell2024whr}. But once we move beyond ordinary human clinical settings, additional questions appear. How should moral weight scale with cognitive or phenomenal capacity? How should one compare individuals to collectives, or biological to artificial systems? What prevents arbitrary double counting when the same localized harms can be described at multiple levels of organization?

The current paper does not attempt a full solution. Instead, it imposes structural constraints: weights should be non-negative, should scale sublinearly with size or complexity, should be approximately invariant under benign coarse-graining, and should not allow the same localized harms to explode through overlapping descriptions. This is less ambitious than a theory of moral weight, but more explicit than simply assuming that all burdens are commensurable. It yields a caring functional that is intended to be empirically interpretable, ethically flexible, and compatible with later specialization.

\subsection{What is new here}

The novelty of the present framework is not a new empirical finding, but a new factorization. Relative to pain science, it distinguishes underlying integrity pressure from low-level valenced channels and from suffering as an architecture-dependent global controller. Relative to metacognitive work, it integrates higher-order access, opacity, and self-indexing into a model of why some aversive states become suffering rather than remaining local control signals. Relative to burden metrics, it provides a pre-aggregation account of what is being aggregated. And relative to discussions of consciousness in this project, it narrows the target from consciousness in general to ethically relevant load-bearing structure, while still using the same architectural vocabulary where useful.

For this reason, the paper should be read less as a finished theory of value than as a coordination object between literatures that usually talk past each other: pain research, metacognition, consciousness theory, and empirical aggregation.


\section{Preliminaries: Agents, Time, and Architecture}

We work in continuous time $t \in \mathbb R$, although nothing essential depends on continuity. Let $\mathcal I$ be a set of \emph{agents} or \emph{entities}. We make no assumption that elements of $\mathcal I$ are all of the same scale: they may be individual organisms, colonies, firms, states, or other Markov-blanketed systems. For each $i \in \mathcal I$ we assume:
\begin{itemize}
  \item a state space $X_i$, in which the agent's relevant internal and boundary states evolve as a function of time;
  \item a collection of control loops indexed by $k \in K_i$, each tracking some state variable $x_{ik}(t)$ against a setpoint or manifold $x^{\ast}_{ik}$;
  \item a bundle of architectural parameters $(B_i,d_i,\beta_i,\tau_i,\dots)$ capturing, in coarse form, bandwidth, recursion depth, self-index strength, opacity, and other structural features.
\end{itemize}

We explicitly \emph{do not} fix a particular cognitive or neural architecture. The role of $(B,d,\beta,\tau)$ is to serve as knobs that modulate how much low-level drive signals show up as what we call suffering.

\section{Integrity Pressure}

\subsection{From homeostatic control to integrity}

Each control loop $k \in K_i$ tracks the deviation of a state variable from its target. Let
\begin{equation}
  e_{ik}(t) = x_{ik}(t) - x^{\ast}_{ik},
\end{equation}
where $x^{\ast}_{ik}$ may be a point, a manifold, or a more complex target set in $X_i$; we suppress this generality in notation. For each loop we posit a non-negative function $h_{ik}$ of the deviation,
\begin{equation}
  h_{ik}: \mathbb R^{n_{ik}} \to [0,\infty), \qquad h_{ik}(0) = 0,
\end{equation}
where $n_{ik}$ is the dimensionality of $e_{ik}$. The function $h_{ik}$ captures how costly a given deviation is for that loop, independent of further architectural amplification.

The \emph{integrity pressure} on agent $i$ at time $t$ is defined as
\begin{equation}
  H_i(t) = \sum_{k \in K_i} w_{ik} \, h_{ik}\bigl(e_{ik}(t)\bigr), \qquad w_{ik} \ge 0. \label{eq:Hi-def}
\end{equation}
Intuitively, $H_i(t)$ measures how many and how severely the agent's homeostatic (and epistemic) controllers are being pushed away from their target regions. Both "bad lack" and "bad excess" of a controlled quantity are handled symmetrically via the same $e_{ik}$; we do not distinguish positive and negative controllers at the level of $H_i$.

\subsection{Black-box semantics and punted details}

Definition~\eqref{eq:Hi-def} is intentionally schematic. We do not specify the microphysical or information-theoretic grounding of $h_{ik}$ or $w_{ik}$, beyond the following commitments:
\begin{itemize}
  \item For each agent $i$, there is some collection of control loops $K_i$ whose deviations matter to its continued existence, organisation, and characteristic functioning.
  \item For each such loop, there is some non-negative mapping $h_{ik}$ that increases (possibly non-linearly and asymmetrically) with the size of the deviation $e_{ik}$ in ways that track the risk of loss of integrity.
  \item The weights $w_{ik}$ encode the relative contribution of each loop to overall integrity pressure.
\end{itemize}

A more detailed treatment could derive $H_i$ from a viability manifold in $X_i$ or from an expected free-energy functional over trajectories. In this paper we only assume that such a derivation is \emph{in principle} possible, and that $H_i(t)$ can be treated as a scalar field over time with the above informal interpretation.

\section{Perceived Pressure, Pain, and Suffering}

\subsection{Perceived integrity pressure}

Agents do not have direct access to $H_i(t)$ as defined in \eqref{eq:Hi-def}. Their controllers operate on a \emph{perceived} integrity pressure $\tilde H_i(t)$ that is implemented in their generative models and control circuitry. In general, $\tilde H_i$ may be biased or distorted relative to $H_i$:
\begin{equation}
  \tilde H_i(t) = F_i\bigl(H_i(t), \xi_i(t)\bigr),
\end{equation}
where $F_i$ encodes both estimation and systematic misrepresentation, and $\xi_i$ captures noise, priors, and structural constraints. Addiction, propaganda, and wireheading are examples where $\tilde H_i$ is driven away from $H_i$.

\subsection{Low-level pain signals}

For each loop $k$ we posit a low-level drive or pain signal $P_{ik}(t)$, which is some function of the local deviation and context:
\begin{equation}
  P_{ik}(t) = f_{ik}\bigl(e_{ik}(t), c_{ik}(t)\bigr),
\end{equation}
where $c_{ik}(t)$ summarises local modulators (e.g. sensitisation, pharmacological state, learned associations). We assume $P_{ik}(t)$ is non-negative and increases with the magnitude of deviations that threaten integrity. Summing over loops yields an aggregate pain signal
\begin{equation}
  P_i(t) = \sum_{k \in K_i} P_{ik}(t).
\end{equation}

This $P_i(t)$ is still pre-phenomenal: it is a bundle of control-relevant signals, not yet what we call suffering.

\subsection{Architectural gain and suffering}

We define \emph{suffering} $S_i(t)$ as the extent to which integrity-relevant deviations are not only registered in local controllers but are also brought into an integrated, self-indexed, temporally extended workspace that can support reportable experience and reflective distress.

We capture this through a gain function $G_i$ that depends on coarse architectural parameters $(B_i,d_i,\beta_i,\tau_i,\dots)$:
\begin{equation}
  S_i(t) \approx P_i(t)\,G_i\bigl(B_i,d_i,\beta_i,\tau_i,\dots\bigr). \label{eq:S-def}
\end{equation}

The architecture parameters discussed here, for illustration and without claiming completeness, are:
\begin{itemize}
  \item $B_i$ is a measure of global bandwidth or workspace capacity, i.e., how many independent signals can be jointly addressed and integrated.
  \item $d_i$ is an effective depth of self-referential recursion: how many layers of "I am aware that I am in pain"-like modelling the system can sustain.
  \item $\beta_i$ encodes the strength and cohesiveness of self-indexing: how sharply the system routes value and control through tokens like "me" or "us".
  \item $\tau_i$ is an opacity parameter, characterising how opaque the underlying processes are to the system's own models: high $\tau$ means the system experiences the outputs of controllers without good access to their generative structure; low $\tau$ means more transparent meta-representation.
\end{itemize}

Equation~\eqref{eq:S-def} is intended as a \emph{local approximation} in ordinary regimes of these parameters. In such regimes it is plausible that increasing $B$, $d$, $\beta$, or $\tau$ tends to increase $S_i$ for fixed $P_i$, at least over some range. Outside those regimes, the relationship may be non-monotonic: for example, extreme increases in $d$ combined with very low $\tau$ might reduce suffering by enabling de-identification with pain, whereas pathological combinations of high $\beta$ and high $\tau$ might amplify it.

The important structural point is that suffering is not a simple readout of $H_i$ or $P_i$. It is the output of a multi-stage process: underlying integrity pressure $H_i$ generates local drive signals $P_{ik}$; those are aggregated into $P_i$; and the architectural gain $G_i$ determines how much of $P_i$ is integrated into a self-indexed, temporally extended phenomenology.

\section{Integrated Integrity Load and the Global Caring Functional}

\subsection{Integrated load per agent}

For each agent $i$ we define its \emph{integrated integrity load}
\begin{equation}
  D_i = \int_{t_0}^{t_1} H_i(t)\,dt,
\end{equation}
where $[t_0,t_1]$ is a time window of interest. We intentionally do not fix a discount function here; $D_i$ is an \emph{undiscounted} quantity. Discounting schemes (for example, exponential or hyperbolic temporal discounting) can be introduced at the level of the global functional if desired, but they are conceptually distinct from what $H_i$ measures.

$D_i$ is the basic object we want to aggregate across agents. It is defined in terms of the underlying $H_i$, not the perceived $\tilde H_i$ or the phenomenological $S_i$. When perceptions are accurate and architectural gain is in a typical human regime, $D_i$ will correlate with total suffering experienced; when perceptions are distorted or architectures differ qualitatively, that correlation can break.

\subsection{Weights and cross-scale aggregation}

We define a global caring functional as
\begin{equation}
  \mathcal L = \sum_{i \in \mathcal I} W_i D_i, \label{eq:L-def}
\end{equation}
where $W_i \ge 0$ is a weight attached to agent $i$. The role of $W_i$ is twofold:
\begin{enumerate}
  \item to capture the intuition that larger, more complex, or more capable minds matter more than smaller or simpler ones, but with diminishing returns;
  \item to handle aggregation across scales in a way that avoids pathological double-counting when we can describe the same underlying situation at multiple levels (individual, group, institution, civilisation, \dots).
\end{enumerate}

Rather than constructing $W_i$ from a specified base level of constituents (cells, neurons, individuals, \dots) or a fixed nesting hierarchy, we treat the set $\{W_i\}_{i \in \mathcal I}$ as a free parameter subject to some broad sanity constraints:
\begin{itemize}
  \item \textbf{Non-negativity}: $W_i \ge 0$ for all $i$.
  \item \textbf{Sublinear scaling}: for agents that intuitively aggregate similar kinds of constituents, the weight of the aggregate grows sublinearly with size or capacity. This expresses "bigger minds matter more, but not in direct proportion to their size".
  \item \textbf{Approximate coarse-graining invariance}: if we simply re-describe the same coherent cluster of constituents at a different scale (e.g. as one village agent or as $N$ individual villagers), the total contribution of that cluster to $\mathcal L$ should not change dramatically.
  \item \textbf{Bounded overcounting}: no localised piece of the world (e.g. a single person or a single colony) should contribute, via overlapping agents, in such a way that its harms are effectively counted an unbounded number of times.
\end{itemize}

Within these constraints, different choices of $W_i$ correspond to different---but structurally similar---ethical positions about how to trade off harms across species and scales. The rest of the framework is intended to be orthogonal to those choices.

\subsection{A reference unit of caring}

Given \eqref{eq:L-def}, we can define a \emph{reference unit of caring} by specifying:
\begin{itemize}
  \item a reference agent $i^\ast$ (for example, a typical human adult in a standard architectural regime);
  \item a reference episode: one second of time during which $H_{i^\ast}(t) = H_0$ for all $t$ in that second, with architectural parameters held at some baseline values $(B_0,d_0,\beta_0,\tau_0,\dots)$;
  \item a reference weight $W_{i^\ast} = 1$.
\end{itemize}

We then define
\begin{equation}
  1 \text{ unit of caring} \equiv \Delta \mathcal L
\end{equation}
associated with preventing that reference episode (i.e. with replacing $H_{i^\ast}(t) = H_0$ by $H_{i^\ast}(t) = 0$ over the specified second, holding everything else fixed). Because $\mathcal L$ is linear in the $D_i$, and $D_{i^\ast}$ is linear in $H_{i^\ast}$, this defines a consistent additive ruler: other interventions can be compared in terms of how many reference-episodes they are equivalent to.

This definition is conventional: different choices of reference agent, episode, or normalisation of $W_{i^\ast}$ yield different numerical units, but the structure of comparisons is preserved.

\section{Relations Between $H$, $P$, $S$, and $\mathcal L$}

The framework separates four layers:
\begin{enumerate}
  \item underlying integrity pressure $H_i(t)$, defined in terms of deviations of control loops that matter for the agent's continued existence and organisation;
  \item perceived integrity pressure $\tilde H_i(t)$, as implemented in the agent's generative models and controllers;
  \item low-level pain or drive signals $P_{ik}(t)$ and their aggregate $P_i(t)$;
  \item phenomenological suffering $S_i(t)$ as in \eqref{eq:S-def}, and the integrated load $D_i = \int H_i(t)dt$ that feeds into $\mathcal L$.
\end{enumerate}

This allows us to express several distinctions:
\begin{itemize}
  \item \textbf{Silent integrity pressure}: high $H_i$ with low $P_i$ and low $S_i$, e.g. when an agent is in grave danger but unaware of it or heavily anaesthetised.
  \item \textbf{Illusory suffering}: high $P_i$ and $S_i$ with low $H_i$, e.g. when perceived threats are badly misaligned with real threats (paranoia, certain phobias, some propaganda and wireheading cases).
  \item \textbf{Architectural modulation}: fixed $H_i$ and $P_i$ with varying $G_i$, e.g. when training, meditation, pharmacology, or pathology change bandwidth, recursion depth, self-indexing, or opacity.
\end{itemize}

In terms of the global caring functional, $\mathcal L$ is tied directly to $H_i$ via $D_i$, not to $S_i$. However, in practice we often use $S_i$ (or proxies thereof) as \emph{evidence} about $H_i$. Inference from $S_i$ to $H_i$ requires empirical models of the mapping from $H_i$ to $P_i$ to $S_i$, including $G_i$ and the distortion operator $F_i$ that produces $\tilde H_i$.

\section{Discussion}

\subsection{Under-specification and degrees of freedom}

The present framework leaves several major choices open:
\begin{itemize}
  \item the microphysical or variational derivation of $H_i$ from dynamics in $X_i$;
  \item the precise functional forms of $h_{ik}$, $f_{ik}$, and $G_i$;
  \item the construction of $W_i$ within the stated constraints;
  \item any temporal discounting or risk-averse deformation applied to $\mathcal L$.
\end{itemize}

This under-specification is deliberate. The aim of the ``unit of caring'' is not to settle normative disputes but to factor out a common scaffold within which different ethical views can be expressed as choices of $W_i$, of mapping from physical states to $H_i$, and of how to trade off $D_i$ across agents.

\subsection{Positive valence and flourishing}

The construction above is asymmetric in the sense that $H_i$ is non-negative and $\mathcal L$ aggregates ``integrity load'' that we would like to reduce. Positive valence and flourishing can be accommodated in at least two ways:
\begin{enumerate}
  \item by treating positive experiences as trajectories in which $H_i$ is low and decreasing, with $S_i$ reflecting the \emph{relief} of resolving integrity threats;
  \item by introducing a separate, structurally analogous quantity for \emph{flourishing pressure} and extending $\mathcal L$ to combine suffering reduction and flourishing promotion.
\end{enumerate}

In this draft we focus on the negative side because pain-driven controllers provide the clearest access to the structure we want to formalise. Nothing in the mathematics requires an exclusively negative ethic.

\subsection{Cross-species and cross-scale comparisons}

For cross-species and cross-scale comparisons, the key questions are:
\begin{itemize}
  \item how $H_i$ scales with the complexity and capabilities of different agents;
  \item how architectural parameters $(B_i,d_i,\beta_i,\tau_i)$ differ and how that affects $G_i$;
  \item how to choose $W_i$ so that different agents' $D_i$ can be compared without pathological double-counting.
\end{itemize}

The present formalism does not answer these questions, but it sharpens them. For example, debates about the moral status of large language models, non-human animals, or entire institutions can be recast as debates about whether there exist plausible mappings from their states to $H_i$, what their architectural parameters look like, and what ranges of $W_i$ would be reasonable under the stated constraints.

\section{Conclusion}

We have introduced an integrity-based notion of harm pressure $H_i(t)$, distinguished it from perceived pressure $\tilde H_i$ and phenomenological suffering $S_i(t)$, and proposed a global caring functional $\mathcal L = \sum_i W_i D_i$ with qualitatively constrained weights $W_i$. This allows us to define a reference unit of caring and to relate it systematically to architectural features of agents and to cross-scale aggregation issues.

Much of the structure is intentionally left schematic: the point is not to close the theory but to make explicit where empirical and normative degrees of freedom enter. The hope is that, by separating integrity pressure, phenomenology, architecture, and aggregation, the framework can serve as a common language for thinking about how much is at stake when we intervene in complex systems.

\bibliographystyle{plainnat}
\bibliography{unit-of-caring}

\end{document}
